% Section fragment. Standard amsmath, amssymb, booktabs, longtable, hyperref.
% Ordinary mathematical proof, with inherited calibration explicitly named.
\section{The PA194QH6 construction and its consistency theorem}
\label{sec:qh6-preservation}

PA194QH6 combines the premise-exchanged distribution emitter of PA194Q
with the coefficient-six constructors and no-push, cached-test controller
of PA194H6. These component ideas are prior work. The new result is this
particular composition, its complete transition table, and the exact source
and physical certificates. It has 194 working states; it does not establish
a reduction below 194.

The theorem is an ordinary mathematical PA-provability statement:
\[
 \mathrm{PA}\vdash
 \bigl(\neg\operatorname{Halt}(\mathrm{PA194QH6},\mathrm{blank})
       \leftrightarrow\operatorname{Con}(\mathrm{PA})\bigr).
 \tag{QH6}
\]
The proof below uses the supplied PA calibration and explicitly shows that
the new source presents an equivalent numerical calculus and admits a
complete proof-to-command compiler. The physical certificates connect that
source to the delivered table. No new end-to-end Lean or Rocq theorem is
claimed. In particular, executing the finite startup is not evidence that
the infinite blank computation never halts.

\subsection{Exact machine and allocation}

The tape alphabet is $\{0,1\}$, with a two-way infinite all-zero initial
tape. The head starts at position zero in dense state zero. Working states
are $0,\ldots,193$ and 194 is the separate halt label. Every transition
writes one bit and moves exactly one cell left or right. The definitive
table is \texttt{machine/machine.tm}; it has 388 rows. Native labels start
at 74 and halt at 254, and \texttt{machine/labels.json} gives the complete
bijection, including the halt label.

The fixed part consists of the same 135 controls as PA194H6. The Q emitter
tree contributes 67 reader controls. Its seven certified history folds
remain valid. One additional fold removes native state 41 in favor of 99,
with final actions
\[
 \delta(99,0)=(0,L,30),\qquad
 \delta(99,1)=(1,L,2040).
\]
All incoming transitions to the removed state leave a zero immediately
to the left of the head; those to the kept state leave a one. The four
possible combinations of this history bit and the current symbol are
checked on the final, simultaneous quotient, not merely on separate
intermediate tables. The same check covers all eight folds, for 32 local
cases. Thus the state count is exactly
\[
                         135+67-8=194.
\]
Every working state is reachable in the directed action graph from the
designated start. This graph property is not a claim that every state is
visited by the actual blank computation.

\subsection{The complete numerical calculus}

Let $p(a,b)=(a+b)(a+b+1)/2+b$ be Cantor pairing. Define
\[
 I(a,b)=6p(a,b),\quad M(a,b)=6p(a,b)+1,\quad
 U(a,B)=6p(B,a)+2,\quad E(a,b)=6p(a,b)+3.
\]
The reversed Cantor arguments in $U$ are intentional. Put
$N(B)=I(B,1)$, $A\wedge B=N(I(A,N(B)))$, and
$\exists_a B=N(U(a,N(B)))$. In formula displays $\to$ means $I$.
Code zero denotes truth, code one is the designated contradiction, and
$I(0,0)=0$ is a real numerical alias.

The six axiom families are
\begin{align*}
 h_1(D,H)&=(\exists_D U(D,H))\to H,\\
 h_2(D,A,H)&=U(D,U(A,H))\to U(A,U(D,H)),\\
 h_3(D,A,H)&=M(H,D)\to(E(H,A)\to M(A,D)),\\
 h_4(D,A,H)&=E(H,D)\to(M(A,H)\to M(A,D)),\\
 h_5(D,H)&=\exists_D E(H,D),\\
 h'_6(D,A,H)&=U(D,A\to H)\to(U(D,A)\to U(D,H)).
\end{align*}
The propositional family is
\[
 \rho_8(D,A,H,u)=((H\to u)\to D)\to((D\to H)\to(A\to H)),
 \qquad u\leq D.
\]
There is a diagonal family $\Gamma(B)=I(B,U(B,B))$. Finally define
\begin{align*}
 V(a,b,c)&=M(a,b)\to U(c,M(a,b)),\\
 P(a,b,c)&=c\to M(c,b),\\
 S(a,b,c)&=U\bigl(b,(M(a,b)\to P(a,b,c))
                         \wedge(M(c,b)\to c)\bigr),\\
 Q(a,b,c)&=U(c,S(a,b,c)\to N(P(a,b,c))),\\
 G(a,b,c)&=N(V(a,b,c)\to Q(a,b,c)).
\end{align*}
The gate $G(a,b,c)$ is an axiom when $c\ne a$ and $c\ne b$.
The only inference rules are modus ponens and generalization
$B\mapsto U(i,B)$ on derivations without undischarged hypotheses. The
formulas may have free labels. Call this numerical calculus
$\mathcal D_6$. A failed gate emission produces zero; zero is already
derivable because $\rho_8(0,0,0,0)=0$.

\subsection{Premise exchange preserves the calculus}

Let $\mathcal C_6$ be the supplied separated-head coefficient-six calculus,
whose sixth axiom is
\[
 h_6(D,A,H)=U(D,A)\to(U(D,A\to H)\to U(D,H)).
\]
All its other families and rules are exactly those of $\mathcal D_6$.
The unchanged restricted $\rho_8$ basis derives the usual implication
axioms $K$ and $S$. The inherited finite derivation in
\texttt{RestrictedImplication.lean} is parametric in implication and uses
only its two argument bounds and reflexivity of $\leq$. Thus it applies
to $I(a,b)=6p(a,b)$ as well as to coefficient four.

From $K,S$ and modus ponens there is the fixed theorem
\[
 (P\to(Q\to R))\to(Q\to(P\to R)).               \tag{E}
\]
The delivered \texttt{proof/exchange\_certificate.py} constructs a
61-line derivation, then independently replays its serialized lines.
It has no hypotheses. Swapping $P,Q$ in the same certificate gives the
opposite direction. Instantiate (E) with
$P=U(D,A)$, $Q=U(D,A\to H)$, $R=U(D,H)$. Each sixth axiom of either
calculus is thereby derived in the other by one axiom, the fixed
certificate, and modus ponens.

Recursion on a finite derivation replaces just its sixth-axiom leaves
by these templates. This yields primitive-recursive maps in both
directions, preserving every conclusion code, whose correctness PA
proves by induction on derivations. No code-order comparison is added
by this exchange. No injectivity of the numeric implication constructor
is required: substituting numbers into a formal $K,S,\mathrm{MP}$
certificate preserves all the defining equalities even when different
expressions denote zero. Consequently PA proves
\[
 \exists d\,\operatorname{Prf}_{\mathcal D_6}(d,1)
 \ \leftrightarrow\
 \exists d\,\operatorname{Prf}_{\mathcal C_6}(d,1).       \tag{1}
\]

\subsection{The inherited PA calibration and its exact scope}

The retained original PA195 calibration, diagonal-calculus proof,
separated-head proof, and \texttt{radix\_calibration.tex} establish the
ordinary mathematical interface
\[
 \mathrm{PA}\vdash
 \bigl(\exists d\,\operatorname{Prf}_{\mathcal C_6}(d,1)
 \leftrightarrow
 \exists p\,\operatorname{Prf}_{\mathrm{PA}}(p,\ulcorner0=1\urcorner)
 \bigr).                                                \tag{2}
\]
This interface, rather than an appeal to a standard model, is the logical
dependency of (QH6). Several details of that interface matter here.

First, recoding from coefficient four to six is hereditary on both
formula codes and variable labels. Define $F(0)=0$ and, for $n>0$,
\[
 n=4p(a,b)+r\ (0\leq r<4)\quad\Longrightarrow\quad
 F(n)=6p(F(a),F(b))+r.
\]
Every decoded component is smaller than $n$, so this is course-of-values
primitive recursion. A decoder $J$ from coefficient six to four has
$J(0)=0$ and, for positive inputs, uses
$J(6p(a,b)+r)=4p(J(a),J(b))+r$ when $r<4$ and returns zero otherwise.
PA proves $J(F(n))=n$ by induction. Hence $F$ is injective and preserves
the disequality guards. It commutes with all four constructors and maps
$\Gamma_4(B)$ to $\Gamma_6(F(B))$ because it also recodes the self-indexed
binder. It fixes zero and one.

Second, $F$ need not preserve numerical order. A restricted
$\rho_{8,4}$ leaf is therefore not translated by assuming
$F(u)\leq F(D)$. The finite \texttt{rho8All} template derives
unrestricted $\rho_8$ from the restricted family at every coefficient
$k\geq4$. Its restricted leaves use only $a\leq a$ and the bounds
$a,b\leq kp(a,b)$. Substituting the recoded parameters in that template
supplies valid target leaves, irrespective of the order of the recoded
numbers. Every offset subtraction in the target command compiler is
subsequently performed on one of those actual target leaves satisfying
its own $u\leq D$ premise.

Third, numbers congruent to four or five modulo six are allowed as
parameters. The total syntactic arithmetic interpretation sends zero
and those unused tags to $0=0$. For a positive number $6p(a,b)+r$, tags
zero through three respectively mean implication, the finite-set bit
predicate $\operatorname{Bit}(x_a,x_b)$, $\forall x_b$ of the first
component, and $x_a=x_b$. Every recursively interpreted component is
smaller than its parent, and every source label occurring in the
interpretation of $B$ is strictly below $B$. Thus $\Gamma_6(B)$ is
uniformly sound by vacuity, including the unused tags. The alias
$I(0,0)=0$ is handled by the fixed equivalence between truth and
truth-implies-truth.

Under this interpretation code one is $\operatorname{Bit}(x_0,x_0)$,
whose negation PA proves. The gate supplier uses finite bit comprehension
for an arbitrary arithmetic formula, with one appropriate PA induction
instance. It is a primitive-recursive supplier of a syntactic PA proof;
it does not decide the formula and invokes no truth or reflection
principle. The witness for a distinct pair of labels has bounding
parameter $x_a+x_b+1$ and characteristic condition
$(\operatorname{Bit}(x_a,z)\vee z=x_b)\wedge\phi[z/x_b]$;
for $a=b$ it uses bound $x_b+1$ and condition
$(z=x_b)\wedge\phi[z/x_b]$. Here $\phi$ interprets $c$,
and the guard together with the code-freshness lemma permits the
new witness at label $c$. These are exactly the two cases retained
in the calibration source.

The original coefficient-four PA-to-calculus direction uses the retained
membership-rewriting, neighborhood-equality, and equality-witness
interfaces. The diagonal and separated-head reports give their uniform
proof replacements. Hereditary recoding and the unrestricted implication
template transfer their finite derivations to coefficient six. In the
other direction, recursive PA proof annotations use the total
interpretation and the schema suppliers just described. These are
quantified, PA-verifiable transformations of proof codes, including
codes in nonstandard models, rather than a separate proof for each
standard input. Equations (1) and (2) calibrate $\mathcal D_6$ to PA.

The numerical stack and scheduler still use ordinary $p$, with no factor
six. Recoding formulas is not asserted to conjugate command stores,
addition, Cantor stacks, or machine trajectories. A new target command
word is compiled from the translated derivation. Replacing the searched
theory by iPRA would also not justify (2): the retained finite-comprehension
supplier uses full PA induction for arbitrary formulas.

\subsection{Exact source and emitter binding}

The input source nodes are exactly those of PA194H6, including
\[
 \mathtt{node163}=(\mathtt{node6},\mathtt{node54}),\qquad
 \mathtt{node170}=((\mathtt{node36},
       (\mathtt{node167},\mathtt{node145})),\mathtt{node57}).
\]
Only the AX grammar is taken from Q. Expanding its macros gives 190
tokens. Relative to the H6 word, only the half-open token interval
$[102,121)$ changes. A free-constructor calculation checks the entire
stack and resulting top for arbitrary symbolic parameters; in particular
the changed interval emits $h'_6$ in place of $h_6$ and does not disturb
the other heads, guard, or parameter registers.

The source binding also checks two exact reconstructions. Using the H6
grammar with the delivered source nodes reproduces the supplied full
H6 instruction graph. Restoring the four-increment multiplier and the
cached-MP node $\mathtt{node163}=(\mathtt{node6},\mathtt{node147})$
reproduces the supplied full Q instruction graph. These are equalities
of all operations and resolved successor targets, not counts or sample
traces.

There are exactly 61 formula-construction occurrences. In every one,
the coefficient-six body is at relative path $1110$ of its constructor,
and there are no other uses of that body. Restoring four increments
gives 3,653 source sites; widening those 61 bodies to six increments
gives 3,775. All 3,409 sites outside the old bodies, and all their
successor targets, agree under an explicit ordinal relocation. No
outside successor enters the interior of a replaced body.

The mathematical loop invariant supplying the changed arithmetic
contract is elementary: after $i\leq Y_0$ iterations of its draining
loop, $T=T_0+6i$ and $Y=Y_0-i$. It therefore terminates with
$(T,Y)=(T_0+6Y_0,0)$ in $8Y_0+1$ primitive source steps. The preceding
Cantor operations are unchanged. For macro coordinates $(T,W,X,Y)$,
with scratch $X=Y=0$, the constructor sends
\[
 (t,p(a,w),0,0)\longmapsto(I(a,t),w,0,0).
\]
The source check binds all formula constructors, including MP matching,
to this loop. The stack-pairing and scheduler-pairing routines do not
use it. The inherited four-increment lemma is not silently cited as
a proof of a six-increment body; the displayed induction is the required
generalization.

\subsection{Command semantics and a complete proof compiler}

A command store is $(t,w,d,h)$: top formula, Cantor stack, numerical
offset, cached formula. Write $[a;w]=p(a,w)$ and use total Cantor
projections, so $\operatorname{hd}0=\operatorname{tl}0=0$ and
$[0;0]=0$. For $b=d+x$ the new controller is
\[
 \mathsf S(t,w,d,h)=
 \begin{cases}\mathsf{halt},&h=1,\\
 (h,\operatorname{tl}w,d,\operatorname{hd}w),&h\ne1,
 \end{cases}
\]
\[
 \mathsf{MP}_x(t,w,d,h)=
 \begin{cases}(b,w,b,0),&I(t,b)=h,\\
 (0,w,b,0),&I(t,b)\ne h.
 \end{cases}                                             \tag{3}
\]
Thus successful MP does not push a dummy zero, and Shift tests only the
cache. Neither feature is new by itself.

For $D=d+x$, $H=h+y$, $A=c$, AX emits the axiom heads of
$\mathcal D_6$, saves the old top, and clears $d,h$:
\[
 \mathsf{AX}_{c,x,y}(t,w,d,h)=
 (G(H,D,A),[q_1,\ldots,q_9,t;w],0,0),
\]
where the gate is replaced by zero on a failed guard,
$q_1=\Gamma(A)$, $q_2=h'_6(D,A,H)$, $q_3=h_5(D,H)$,
$q_4=h_4(D,A,H)$, $q_5=h_3(D,A,H)$, $q_6=h_2(D,A,H)$,
$q_7=h_1(D,H)$, $q_8=\rho_8(D,A,H,d)$, and $q_9=U(d,h)$.
The side condition for $q_8$ holds because $D=d+x$.
All of $q_1,\ldots,q_8$ and the emitted top have residue zero;
$q_9$ has residue two and is at least two. None is one.

Let $A_0=\mathsf{AX}_{0,0,0}$, and execute command words left to right.
The exact identities
\begin{align*}
 \mathsf{AX}_{c,x,y};\mathsf S^9(t,w,d,h)
   &=(q_8,[t;w],0,U(d,h)),\\
 \mathsf{AX}_{c,x,y};\mathsf S^{10}(t,w,d,h)
   &=(U(d,h),w,0,t)
\end{align*}
do not halt, because they test only the cleared cache and the non-one
generated heads. Define
\[
 \mathsf{Push}=A_0;\mathsf S^9;\mathsf{MP}_0,\quad
 \mathsf{Pop}=A_0;\mathsf S^{11},\quad
 \mathsf{Discard}=\mathsf{MP}_0,\quad
 \mathsf{Offset}_i=\mathsf{MP}_i.
\]
Then
\begin{align*}
 \mathsf{Push}(t,w,d,h)&=(0,[t;w],0,0),\\
 \mathsf{Pop}(t,[a;w],d,h)&=
    \begin{cases}\mathsf{halt},&t=1,\\(t,w,0,a),&t\ne1,\end{cases}\\
 \mathsf{Discard}(t,w,0,h)&=(0,w,0,0),\\
 \mathsf{Offset}_i(0,w,0,0)&=(0,w,i,0).
\end{align*}
Push uses the residue mismatch between implication and universal codes.
For Discard the two MP branches coincide when the proposed conclusion is
zero, even with arbitrary old cache. Pop's last Shift tests exactly the
saved top. These are numerical identities and remain valid for a collapsed
zero push or an empty stack. The symbolic command checker explicitly
includes those cases.

From $(t,w,0,0)$ the word
\[
 \mathsf{Discard};\mathsf{AX}_{A,D,H};\mathsf S^{j+1};
                 \mathsf{Pop}^{9-j}\quad(1\leq j\leq7)
\]
selects $q_j$, preserving $w$ and clearing the caches. The gate uses
$\mathsf{Discard};\mathsf{AX}_{c,b,a};\mathsf{Pop}^{10}$.
For a restricted propositional leaf with $u\leq D$, its word is
$\mathsf{Discard};\mathsf{Offset}_u;
\mathsf{AX}_{A,D-u,H};\mathsf S^9;\mathsf{Pop}$.
Generalization is implemented by
\[
 \mathsf{Gen}_i=\mathsf{Push};\mathsf{Offset}_i;\mathsf S;
              \mathsf{AX}_{i,0,0};\mathsf S^{10},
\]
which sends $(a,w,0,0)$ to $(U(i,a),w,0,0)$.

The MP suffix $\mathsf S;\mathsf S;\mathsf{MP}_b;\mathsf{Pop}$
acts on $(0,[a,I(a,b),0;w],0,0)$. It returns $(b,w,0,0)$
when $a,b\ne1$, and halts when $a=1$ or $b=1$.
The zero below the premises is essential: it replaces the dummy zero
that older MP controllers themselves pushed. For a proof tree $p$,
compile its axioms and generalization using these words. At an MP node
with proofs $p_0$ of $a$ and $p_1$ of $I(a,b)$, use
\[
 W_p=\mathsf{Discard};\mathsf{Push};W_{p_1};\mathsf{Push};
 W_{p_0};\mathsf{Push};\mathsf S;\mathsf S;\mathsf{MP}_b;
 \mathsf{Pop}.
\]
Induction on $p$ proves that execution from every $(t,w,0,0)$ either
halts or returns $(\operatorname{concl}p,w,0,0)$. In the MP case
the first two commands reserve zero before the two subproofs.
Therefore if $p$ proves one, $A_0;W_p;\mathsf{Pop}$ halts from
\emph{every} incoming store. The initial AX clears arbitrary carried
caches; the final Pop tests one before using the old underlying tail.
A final plain Shift would not be correct, since
$\mathsf S(1,0,0,0)=(0,0,0,0)$.

\subsection{Scheduler, soundness, and the literal theorem}

The unchanged scheduler enumerates command words while retaining the
current store between trials. The all-store property just proved is
therefore necessary. The continuation encoding is $e(\epsilon)=0$ and
$e(c::w)=1+b_c(e(w))$, where
\[
 b_{\mathsf S}(z)=p(0,z),\quad
 b_{\mathsf{MP}_x}(z)=p(1,p(x,z)),\quad
 b_{\mathsf{AX}_{c,x,y}}(z)=p(c+2,p(y,p(x,z))).
\]
If $z=e(A_0;W_p;\mathsf{Pop})$, prefixing this halting word by
$\mathsf{AX}_{k,0,0}$ gives the increasing, unbounded candidate codes
$p(k+2,p(0,p(0,z)))$. A prefix can alter the store without harm because
the suffix halts from every store. Fresh candidates begin at two.
Within a candidate every positive continuation strictly decreases,
and the arithmetic source loops have primitive-recursive termination
bounds. PA verifies by bounded candidate induction that a member of
this cofinal family is reached unless an earlier halt occurs.

For the converse, annotate $t,h$ and every canonical Cantor-stack entry
by a derivation in $\mathcal D_6$. The tail of a positive Cantor stack
is strictly smaller, so its canonical finite list and annotations are
primitive-recursively coded. Virtual zero entries carry the fixed proof
of zero. AX preserves annotations by its generated axioms and
generalization; successful MP combines its two annotated premises;
failed MP gives zero; Shift moves annotations and halts only when an
annotated cache equals one. Thus every finite command halt yields a
$\mathcal D_6$ refutation. PA proves this invariant by induction and
then applies (1)--(2) to obtain a PA refutation.

The delivered universal physical certificates check all 4,876 outcomes
of the 3,775-site source. They quantify over all nine natural primary
registers, both decrement branches, header gap $48+n-|p|$ at source
address $p$ for every $n\geq0$,
and any nonempty finite list of natural spare registers. The independent
checker derives its four register-list lemmas from the literal rows,
uses two independent source resolvers, and verifies the dense/native
bijection. Its 1,169,239 symbolic operations have maximum case length
403 and minimum positive physical-duration constant three. The separate
construction checker uses 1,601,532 symbolic operations, with maximum
499. All eight folds are rechecked on the simultaneous final table.

Exactly 58,088 literal steps from blank reach the clean source entry
25 at native state 84 and head $-2$, with logical registers
$(2,0,0,0,0,0,0,0,0)$ and 26 unit spare blocks. This check compares
the complete finite-support tape, head, control and zero exterior.
The finite row equalities, unary-run inductions, register-list
inductions, source-loop invariants, and finite-trace concatenations
are formalizable in PA; the checks are executable certificates for
the specific finite bindings, rather than newly kernel-checked proofs.

For a PA refutation, (1)--(2) give a numerical refutation, the command
compiler gives an all-store halting word, the scheduler reaches it or
halts earlier, and the physical contracts give a finite literal blank
halt. Conversely, under $\operatorname{Con}(\mathrm{PA})$ the annotated
command invariant precludes any command halt. For a proposed literal
halting time $T$, iterate $T+1$ positive completed source segments and
their physical simulations. After adjoining startup their exact run
extends beyond $T$. Determinism contradicts the alleged earlier halt.
This argument needs neither inverse decoding of arbitrary intermediate
tapes nor unbounded search for the next boundary. All operations and
correctness inductions are on finite codes in PA. The two directions
establish (QH6), with the explicitly stated inherited calibration.

\subsection{Material difference and limits}

Against H6, QH6 changes the recurring emitted distribution formula
$P\to(Q\to R)$ to $Q\to(P\to R)$, while preserving the surrounding
axiom stack and maintaining a complete proof compiler. Against Q and
QN, it also uses coefficient-six formula coding; against QN its Shift
tests only $h=1$. These change recurring command semantics. The result
is a new composition of known ideas, not a claim that premise exchange,
coefficient six, or either controller policy was invented here.

An exact action-graph census of the pinned repository, including its
current incoming archives, finds no prior table isomorphic to this one,
even permitting arbitrary state renaming, any spanning start, reflection,
and a global exchange of symbols. That finite novelty statement concerns
the recognized literal-table corpus and does not establish inequivalence
of arbitrary algorithms, global historical priority, or minimality of
194 states.
